% ======================================================================
% Appendix B: The convolutional PML in the Hamiltonian-simulation
% variables.  Self-contained fragment, \input from main.tex after the
% \appendix switch.  All labels are prefixed appB:.
% ======================================================================
\appendix{The convolutional PML in the Hamiltonian-simulation variables}
\label{ch3:appB:cpml}

This appendix expands the one-line statements of
Sections~\ref{ch3:sec:cpml}, \ref{ch3:sec:1d}, and \ref{ch3:sec:2d}: the causal kernel
and memory ODE from the CFS stretch
(Appendix~\ref{ch3:appB:kernel}), the one-dimensional memory generator and
the frequency-domain collapse (Appendix~\ref{ch3:appB:1d}), the proof that
the collapse is exact (Appendix~\ref{ch3:appB:collapse}), the
two-dimensional eight-block generator
(Appendix~\ref{ch3:appB:2d}), the staggered profile-sampling rule
(Appendix~\ref{ch3:appB:stagger}), and the profile formulas
(Appendix~\ref{ch3:appB:profiles}).  Throughout, $D$ is the
forward-difference matrix and $D^{\dagger}=-D_{\mathrm{b}}$ its negative
backward counterpart (Section~\ref{ch3:sec:wave}); all operators are exactly
those constructed in the accompanying code.

\subsection{From the CFS stretch to the memory ODE}
\label{ch3:appB:kernel}

We fix the Fourier convention
$\hat f(\omega)=\int dt\,e^{-i\omega t}f(t)$,
$f(t)=(2\pi)^{-1}\int d\omega\,e^{+i\omega t}\hat f(\omega)$, so that
$\partial_t\mapsto i\omega$ and the causal exponential has the transform
$\int_0^{\infty}dt\,e^{-\beta t}e^{-i\omega t}=(\beta+i\omega)^{-1}$ for
$\beta>0$.  A PML stretches the coordinate into the complex plane,
$\tilde x=\int_0^x s(x',\omega)\,dx'$, which in the frequency domain
replaces every spatial derivative by
$\partial_x\mapsto s(\omega)^{-1}\partial_x$ inside the layer.  For the
CFS factor of Eq.~\eqref{ch3:eq:cfs}, partial fractions in the variable
$i\omega$ give
\begin{equation}
  \frac{1}{s(\omega)}
  \;=\; \frac{\alpha+i\omega}{\kappa(\alpha+i\omega)+\sigma}
  \;=\; \frac{1}{\kappa} \;+\; \hat\zeta(\omega),
  \label{ch3:appB:eq:invs}
\end{equation}
with
\begin{equation}
  \hat\zeta(\omega) \;=\; -\frac{\sigma}{\kappa^{2}}\,
  \frac{1}{\beta+i\omega},
  \qquad
  \beta \;:=\; \frac{\sigma}{\kappa}+\alpha
  \label{ch3:appB:eq:zetahat}
\end{equation}
the CFS decay rate: indeed
$1/s-1/\kappa=(\kappa-s)/(\kappa s)$ with
$\kappa-s=-\sigma/(\alpha+i\omega)$ and
$\kappa s=\kappa^{2}(\beta+i\omega)/(\alpha+i\omega)$.  The single pole
at $i\omega=-\beta$ inverts, by the causal transform above, to the
exponential kernel of Eq.~\eqref{ch3:eq:kernel},
$\zeta(t)=-(\sigma/\kappa^{2})\,e^{-\beta t}\theta(t)$, and frequency
multiplication is time convolution:
$s^{-1}\partial_x f = \kappa^{-1}\partial_x f + \zeta * \partial_x f$.

The convolution becomes local in time once it is named.  Let
$\phi:=\zeta*\partial_x f$.  Multiplying Eq.~\eqref{ch3:appB:eq:invs} by
$(\beta+i\omega)$ gives the polynomial identity
$(\beta+i\omega)\hat\zeta=-\sigma/\kappa^{2}$, hence
$i\omega\,\hat\phi=-\beta\,\hat\phi
-(\sigma/\kappa^{2})\,\widehat{\partial_x f}$, which transforms back to
the memory ODE of Eq.~\eqref{ch3:eq:memoryode}.  Equivalently, in the time
domain,
$\partial_t\phi=\zeta(0^{+})\,\partial_x f+\zeta'*\partial_x f$ with
$\zeta(0^{+})=-\sigma/\kappa^{2}$ and $\zeta'=-\beta\zeta$.  One
first-order ODE per stretched derivative thus carries the convolution
exactly; this is the recursive-convolution CPML of Roden and
Gedney~\cite{roden2000cpml}.

\subsection{One dimension: memory generator and collapse}
\label{ch3:appB:1d}

In the variables of Section~\ref{ch3:sec:wave} the undamped 1D system is
$\dot v=-iDw$, $\dot w=-iD^{\dagger}v$ [Eq.~\eqref{ch3:eq:waveH}].  Each
equation contains exactly one spatial derivative, so the stretch
generates one memory field per physical field.  Sampling the profiles
per field (Appendix~\ref{ch3:appB:stagger}) defines the diagonal matrices
$\Sigma_v=\diag(\sigma_{v,j})$, $K_v=\diag(\kappa_{v,j})$,
$B_v=\diag(\sigma_{v,j}/\kappa_{v,j}+\alpha_{v,j})$, and the layer
projector $\Pi_v=\diag(\mathbf{1}[\sigma_{v,j}>0])$, and likewise for
$w$.  Applying the stretched-derivative rule to $-iDw$ and carrying the
convolution by the rescaled memory field
$\phi_v:=\gamma^{-1}\,\zeta_v*(-iDw)$---where $\gamma>0$ is a free
constant rescaling under which the physical trajectory is invariant,
set to the balanced value $\gamma=\sqrt{2\sigma_{\max}/h}$ of
Appendix~\ref{ch3:app:proof}---gives, by Eq.~\eqref{ch3:eq:memoryode},
\begin{equation}
  \dot v = -iK_v^{-1}Dw + \gamma\phi_v,
  \qquad
  \dot\phi_v = -B_v\phi_v + \frac{i}{\gamma}\Sigma_vK_v^{-2}Dw,
  \label{ch3:appB:eq:stretched}
\end{equation}
and the analogous pair for $(w,\phi_w)$ with $D\to D^{\dagger}$,
$v\leftrightarrow w$.  The source of $\phi_v$ carries the factor
$\Sigma_v$, so $\phi_v$ vanishes identically outside the layer strips
for memory-free data; masking the injection column by $\Pi_v$ therefore
changes no such trajectory while removing the spurious $O(\gamma)$
interior entries that an unmasked column would contribute to the
Hermitian part $H_1$.  On the four-block state
$[v;\,w;\,\phi_v;\,\phi_w]$ (two extra block qubits) the generator is
\begin{equation}
  A_{\mathrm{mem}} =
  \begin{bmatrix}
    0 & -iK_v^{-1}D & \gamma\Pi_v & 0\\[2pt]
    -iK_w^{-1}D^{\dagger} & 0 & 0 & \gamma\Pi_w\\[2pt]
    0 & \tfrac{i}{\gamma}\Sigma_vK_v^{-2}D & -B_v & 0\\[2pt]
    \tfrac{i}{\gamma}\Sigma_wK_w^{-2}D^{\dagger} & 0 & 0 & -B_w
  \end{bmatrix}.
  \label{ch3:appB:eq:gen1d}
\end{equation}

\emph{Collapse for $\kappa=1$, $\alpha=0$.}  In this case the stretch
degenerates: $s(\omega)=(i\omega+\sigma)/(i\omega)$, so the stretched
$v$ equation, $i\omega\hat v=s_v(\omega)^{-1}(-iD\hat w)$ per Fourier
mode, multiplies out to
\begin{equation}
  (i\omega+\sigma_v)\,\hat v \;=\; -iD\hat w,
  \label{ch3:appB:eq:freqcollapse}
\end{equation}
because the symbol $s(\omega)\,i\omega=i\omega+\sigma$ is a
\emph{first-degree polynomial} in $i\omega$.  Transforming back,
$\dot v=-iDw-\sigma_v v$, and likewise
$\dot w=-iD^{\dagger}v-\sigma_w w$: the convolution collapses to the
local damping of Eq.~\eqref{ch3:eq:1dcollapse},
$A=-iH-\diag(\sigma_v,\sigma_w)$, with no memory fields at all.

\subsection{Proof of the exact collapse}
\label{ch3:appB:collapse}

The frequency-domain argument above is formal (it commutes the stretch
with the time-domain initial-value problem); the following statement
makes the equivalence exact at the semi-discrete level.

\begin{lemma}[exact collapse]
\label{ch3:appB:lem:collapse}
Let $\kappa\equiv1$ and $\alpha\equiv0$ in the generator of
Eq.~\eqref{ch3:appB:eq:gen1d}, so that $B_v=\Sigma_v$ and
$B_w=\Sigma_w$, and define the defects
\begin{equation}
  e_v := \phi_v + \gamma^{-1}\Sigma_v v,
  \quad
  e_w := \phi_w + \gamma^{-1}\Sigma_w w.
  \label{ch3:appB:eq:defect}
\end{equation}
Then $\dot e_v=\dot e_w=0$ along every solution of
$\dot\Psi=A_{\mathrm{mem}}\Psi$.  Consequently a memory-form trajectory
projects onto a trajectory of the collapsed generator
Eq.~\eqref{ch3:eq:1dcollapse} if and only if $e_v(0)=e_w(0)=0$; for the
standard embedding $\phi_v(0)=\phi_w(0)=0$ this is the condition
$\Sigma_v v(0)=0$ and $\Sigma_w w(0)=0$, i.e.\ initial data supported
outside the layer.
\end{lemma}

\begin{lemproof}
Using both rows of Eq.~\eqref{ch3:appB:eq:stretched} with $B_v=\Sigma_v$,
\begin{equation*}
\begin{split}
  \dot e_v
  &= \Bigl[-\Sigma_v\phi_v+\tfrac{i}{\gamma}\Sigma_vDw\Bigr]
  + \gamma^{-1}\Sigma_v\bigl[-iDw+\gamma\Pi_v\phi_v\bigr]\\
  &= (\Sigma_v\Pi_v-\Sigma_v)\,\phi_v \;=\; 0,
\end{split}
\end{equation*}
since $\Sigma_v\Pi_v=\Sigma_v$ (the projector is the support indicator
of $\Sigma_v$); identically for $e_w$.  Equivalently, the block-row
operator
$E=\bigl[\begin{smallmatrix}
\gamma^{-1}\Sigma_v & 0 & I & 0\\
0 & \gamma^{-1}\Sigma_w & 0 & I
\end{smallmatrix}\bigr]$
annihilates the generator from the left, $EA_{\mathrm{mem}}=0$
(verified to round-off in our implementation).  On the invariant slice
$e\equiv0$ the memory ansatz $\phi_v=-\gamma^{-1}\Sigma_v v$ of
Section~\ref{ch3:sec:1d} holds at all times, the injection term becomes
$\gamma\Pi_v\phi_v=-\Sigma_v v$, and $(v,w)$ obeys
Eq.~\eqref{ch3:eq:1dcollapse}; conversely, every collapsed trajectory lifts
to the memory form through the ansatz.
\end{lemproof}

When $e(0)\neq0$, the conserved defect acts on the physical equations
as the static in-layer forcing $\gamma\Pi_ve_v(0)$,
$\gamma\Pi_we_w(0)$, so the discrepancy between the two forms scales
with the layer content $\lVert\Sigma_vv(0)\rVert$,
$\lVert\Sigma_ww(0)\rVert$ of the initial data---and vanishes exactly
for data supported outside the layer.  Numerically, propagating the
same interior pulse with Eq.~\eqref{ch3:appB:eq:gen1d} and with
Eq.~\eqref{ch3:eq:1dcollapse} agrees to $6\times10^{-13}$
(Section~\ref{ch3:sec:1d}).

\emph{Remark (the roles of $\alpha$ and $\kappa$).}  The cancellation
uses only $\alpha\equiv0$: for a graded $\kappa$ the modified defect
$e_v=\phi_v+\gamma^{-1}\Sigma_vK_v^{-1}v$ is conserved by the same
computation, and the convolution collapses to the local form
$\dot v=-iK_v^{-1}Dw-K_v^{-1}\Sigma_vv$ (both statements again verified
to round-off).  For $\alpha>0$, instead, $\dot e_v=-\diag(\alpha_v)\phi_v
\neq0$---equivalently, $s(\omega)\,i\omega$ is no longer polynomial in
$i\omega$---and the memory fields are irreducible: the four-block form
is then the genuine content of the CFS layer.

\subsection{Two dimensions: the eight-block generator}
\label{ch3:appB:2d}

In 2D the damping acts per direction: the $x$ stretch replaces every
$\partial_x$ and the $y$ stretch every $\partial_y$.  On the
separated-block state of Section~\ref{ch3:sec:wave} the undamped equations are
$\dot v=-iD_xw_x-iD_yw_y$, $\dot w_x=-iD_x^{\dagger}v$,
$\dot w_y=-iD_y^{\dagger}v$, with $D_x=D\otimes I_{N_y}$ and
$D_y=I_{N_x}\otimes D$ on the flattened $N_xN_y$ grid register ($x$
register more significant).  The $v$ equation contains one derivative
of each direction and acquires two memory fields $\phi_{vx}$,
$\phi_{vy}$; each $w$ equation contains a single derivative and
acquires one, $\phi_x$ for $w_x$ and $\phi_y$ for $w_y$.  The four
memory fields plus the three physical fields are padded by one zero
block to the eight blocks of Eq.~\eqref{ch3:eq:blocks} (three block
qubits).

Because $\sigma_x$ depends on $x$ only, its diagonal samples factorise
on the grid register:
$\Sigma_v^{x}=S_x\otimes I_{N_y}$ and
$\Sigma_w^{x}=S_x^{\mathrm{st}}\otimes I_{N_y}$, where
$S_x=\diag\,\sigma(x_j)$ holds the nodal samples and
$S_x^{\mathrm{st}}=\diag\,\sigma(x_{j-1/2})$ the half-cell samples
(Appendix~\ref{ch3:appB:stagger}); symmetrically
$\Sigma_v^{y}=I_{N_x}\otimes S_y$ and
$\Sigma_w^{y}=I_{N_x}\otimes S_y^{\mathrm{st}}$.  With $\kappa\equiv1$
(our 2D configurations), the decay rates are
$B_v^{x}=\Sigma_v^{x}+\diag\alpha_x(x_j)$ etc., and the masks
$\Pi_v^{x}=\diag(\mathbf{1}[\Sigma_v^{x}>0])$, \dots\ are the
layer-strip projectors.  Repeating the derivation of
Appendix~\ref{ch3:appB:1d} once per (field, stretched direction) pair gives
the generator, in the block ordering of Eq.~\eqref{ch3:eq:blocks},
\begin{widetext}
\begin{equation}
  A \;=\;
  \begin{bmatrix}
    0 & -iD_x & -iD_y & \gamma\Pi_v^{x} & \gamma\Pi_v^{y} & 0 & 0 & 0\\[3pt]
    -iD_x^{\dagger} & 0 & 0 & 0 & 0 & \gamma\Pi_w^{x} & 0 & 0\\[3pt]
    -iD_y^{\dagger} & 0 & 0 & 0 & 0 & 0 & \gamma\Pi_w^{y} & 0\\[3pt]
    0 & \frac{i}{\gamma}\Sigma_v^{x}D_x & 0 & -B_v^{x} & 0 & 0 & 0 & 0\\[3pt]
    0 & 0 & \frac{i}{\gamma}\Sigma_v^{y}D_y & 0 & -B_v^{y} & 0 & 0 & 0\\[3pt]
    \frac{i}{\gamma}\Sigma_w^{x}D_x^{\dagger} & 0 & 0 & 0 & 0 & -B_w^{x} & 0 & 0\\[3pt]
    \frac{i}{\gamma}\Sigma_w^{y}D_y^{\dagger} & 0 & 0 & 0 & 0 & 0 & -B_w^{y} & 0\\[3pt]
    0 & 0 & 0 & 0 & 0 & 0 & 0 & 0
  \end{bmatrix}
  \quad\text{on}\quad
  \bigl[\,v;\; w_x;\; w_y;\; \phi_{vx};\; \phi_{vy};\;
        \phi_{x};\; \phi_{y};\; 0\,\bigr],
  \label{ch3:appB:eq:gen2d}
\end{equation}
\end{widetext}
which is, entry for entry, the matrix assembled by the accompanying
implementation.

\emph{Strip masking.}  The source of each memory field carries its own
profile factor ($\Sigma_v^{x}$ for $\phi_{vx}$, \dots), so the field
remains identically zero outside its layer strip for memory-free
initial data; the corresponding injection column is masked by the strip
projector, exactly as in 1D, so that $H_1$ carries no spurious
interior coupling.

\emph{Corners.}  In a corner cell both $\sigma_x>0$ and $\sigma_y>0$
hold simultaneously: both directions' memory fields are active there,
the same generator entries of Eq.~\eqref{ch3:appB:eq:gen2d} apply
unchanged, and no corner-specific terms arise.  This is the structural
advantage of the unsplit convolutional PML over split-field
formulations~\cite{roden2000cpml}.  The measured stability,
$\max\operatorname{Re}\operatorname{eig}(A)=0$ to round-off, is
reported in Section~\ref{ch3:sec:2d}.

\subsection{The staggered sampling rule}
\label{ch3:appB:stagger}

For the forward-difference convention, $(Dw)_j=(w_{j+1}-w_j)/h$.  The
component $v_j$ is updated through $(Dw)_j$ and represents the field at
the integer node $x_j=jh$; the component $w_j$ is updated through
$(D^{\dagger}v)_j\propto v_j-v_{j-1}$, a difference centered at the
half-cell point $x_{j-1/2}$, so the $w$ field lives on the staggered
grid (the backward convention shifts it to $x_{j+1/2}$).  The sampling
rule follows: \emph{every profile factor in an evolution equation is
sampled on the grid of the field that the equation updates},
\begin{equation}
  \sigma_{v,j}=\sigma(x_j),
  \qquad
  \sigma_{w,j}=\sigma(x_{j\mp1/2}),
  \label{ch3:appB:eq:staggerrule}
\end{equation}
with the upper (lower) sign for the forward (backward) convention,
and per direction in 2D (\ref{ch3:appB:tab:stagger}): only the $x$
profile of the $w_x$ family and the $y$ profile of the $w_y$ family are
staggered; all other dependencies are nodal.

\begin{table}
\caption{Staggered sampling of the damping profiles in 2D
(forward-difference convention; the backward convention flips the
half-cell shifts).  Each profile is sampled on the grid of the field
its equation updates.}
\label{ch3:appB:tab:stagger}
\begin{tabular}{lll}
\toprule
Field & Lives at & Profile samples \\
\midrule
$v$   & $(x_j,\,y_k)$         & $\sigma(x_j)$,\ $\sigma(y_k)$ \\
$w_x$ & $(x_{j-1/2},\,y_k)$   & $\sigma(x_{j-1/2})$ \\
$w_y$ & $(x_j,\,y_{k-1/2})$   & $\sigma(y_{k-1/2})$ \\
\bottomrule
\end{tabular}
\end{table}

\emph{Ablation.}  Sampling all profiles at the integer nodes
instead---the only change being the half-cell shift of the $w$
profiles---degrades the measured gold-standard reflection of
Section~\ref{ch3:sec:2d} by factors of 120--235.  At the reflection floors of
\ref{ch3:fig:convergence}(a,b) the half-cell stagger is load-bearing,
not a refinement.

\subsection{Profile formulas}
\label{ch3:appB:profiles}

Within a layer of $n_{\mathrm{pml}}$ points (width
$L=n_{\mathrm{pml}}h$) the damping is polynomially graded in the depth
$d$ measured from the interface,
\begin{equation}
  \sigma(d) \;=\; \sigma_{\max}\,(d/L)^{m},
  \qquad 0\le d\le L,
  \label{ch3:appB:eq:sigma}
\end{equation}
and zero outside; the interface node itself is undamped,
$\sigma(0)=0$, the outermost node attains $\sigma_{\max}$, and
staggered samples beyond the wall are clipped to $d=L$.  The amplitude
follows from the normal-incidence round trip: a wave crossing the
layer, reflecting off the terminating hard wall, and crossing back
accumulates the attenuation
\begin{equation}
  R \;=\; \exp\Bigl[-\frac{2}{c}\int_0^{L}\sigma(x)\,dx\Bigr]
  \;=\; \exp\Bigl[-\frac{2\,\sigma_{\max}L}{(m+1)\,c}\Bigr],
  \label{ch3:appB:eq:roundtrip}
\end{equation}
so prescribing a design reflection $R_0$ gives the Collino--Tsogka
amplitude~\cite{collino2001pml}
\begin{equation}
  \sigma_{\max} \;=\; -\frac{(m+1)\,c\,\ln R_0}{2L},
  \label{ch3:appB:eq:sigmamax}
\end{equation}
used with the standard recipe $m=2$,
$R_0=10^{-3}$~\cite{komatitsch2007unsplit}.  The nominal $R_0$ is
realised only for layers of ${\gtrsim}8$ points; for thinner layers the
discrete transition error dominates (a four-point layer reflects at the
$0.5$--$2\%$ level for grid-scale wavelengths), consistent with the
saturation of \ref{ch3:fig:convergence}(a).

When the CFS parameters are used they are graded with the same
geometry,
\begin{equation}
  \kappa(d) = 1 + (\kappa_{\max}-1)\,(d/L)^{m},
  \quad
  \alpha(d) = \alpha_{\max}\,(1-d/L),
  \label{ch3:appB:eq:grading}
\end{equation}
$\kappa$ rising from $1$ at the interface to $\kappa_{\max}$ at the
wall with the same polynomial as $\sigma$, and $\alpha$ tapering
linearly~\cite{komatitsch2007unsplit} from $\alpha_{\max}$ just inside
the interface to zero at the wall (and vanishing identically outside
the layer); in the implementation
$\alpha(d)=\alpha_{\max}[1-(\sigma(d)/\sigma_{\max})^{1/m}]$, which is
identical.  A nonzero $\alpha$ improves grazing-incidence and
evanescent absorption but unmatches the layer at low frequencies (the
stretch is matched only for $\omega\gg\alpha$), so the fixed-horizon
absorption benchmarks of Section~\ref{ch3:sec:2d} use $\alpha_{\max}=0$.
Finally, the memory rescaling defaults to
$\gamma=\sqrt{2\sigma_{\max}/h}$, the choice that balances the
injection and memory-source contributions to $\lamax(H_1)$
(Appendix~\ref{ch3:app:proof}).
